\documentclass[11pt]{article}

\usepackage[utf8]{inputenc}
\usepackage[T1]{fontenc}
\usepackage{lmodern}
\usepackage{geometry}
\usepackage{amsmath,amssymb,amsfonts,amsthm,mathtools}
\usepackage{bm}
\usepackage{enumitem}
\usepackage{microtype}
\usepackage{hyperref}
\usepackage{titlesec}
\usepackage{booktabs}
\usepackage{array}
\usepackage{setspace}
\usepackage{mathrsfs}
\usepackage{physics}

\geometry{margin=1in}
\hypersetup{colorlinks=true, linkcolor=black, urlcolor=blue, citecolor=black}
\setlist[enumerate]{topsep=0.4em,itemsep=0.35em}
\setlist[itemize]{topsep=0.3em,itemsep=0.25em}
\titleformat{\section}{\large\bfseries}{\thesection.}{0.5em}{}
\titleformat{\subsection}{\normalsize\bfseries}{\thesubsection.}{0.5em}{}
\setstretch{1.06}

\newcommand{\Aether}{\AE{}ther}
\newcommand{\Aflow}{\AE{}ther-flow}
\newcommand{\TheoryName}{\Aflow{} Interpretation of Relativity}
\newcommand{\TheoryProgram}{\Aether{} / \Aflow{} framework}
\newcommand{\CompletedMetric}{g^{\sharp}}
\newcommand{\CompletionCore}{\mathfrak C^{\mathrm{zr,aff}}}
\newcommand{\TraceCore}{\mathfrak C^{\mathrm{tr}}}
\newcommand{\AffineNucleus}{\mathfrak N^{\mathrm{aff}}}
\newcommand{\AmbientPackage}{\mathfrak A^{\mathrm{bal}}}

\newtheorem{definition}{Definition}
\newtheorem{proposition}{Proposition}
\newtheorem{remark}{Remark}
\newtheorem{corollary}{Corollary}
\newtheorem{theorem}{Theorem}

\title{\textbf{The \TheoryName{}}\\[0.4em]
\large Primitive-Reservoir Observer-Reduction\\
\large Completion-Balance Ambient Dynamics\\
\large Next Benchmark Criterion}
\author{}
\date{}

\begin{document}

\maketitle

\begin{abstract}
The later trace-core collapse theorem sharpened the live internal burden below the completion-core frontier to the smaller completion-state shell-trace core \(\TraceCore\). The new ambient trace-core forcing theorem has now spent the preferred honest move available there: the trace core itself is a canonical and necessary reduction of the deeper completion-balance ambient dynamics package \(\AmbientPackage\). The remaining honest burden is therefore no longer to restate the trace core, the completion core, or the exact-shell affine nucleus.

This note records the routing consequence of that theorem. The next honest benchmark-facing manuscript must therefore either derive or uniquely force the completion-balance ambient dynamics package \(\AmbientPackage\) from still more primitive explicit substrate structure in a way that changes benchmark-facing status, or prove a genuinely smaller theorem-level collapse, sharper necessity result, or sharper uniqueness result strictly inside that ambient package. Another manuscript that merely replays the same ambient trace-core consequence without removing ambient-balance freedom does not count next.

The benchmark boundary remains fixed throughout. The one operative metric is still exactly \(\CompletedMetric\), the ordinary matter route is unchanged, there is no second operative metric, the low-energy matter law is not enlarged, and stronger first-principles promotion language remains paused. The positive result of this note is therefore routing discipline at the new ambient-dynamics frontier.
\end{abstract}

\section{Introduction}

The active derivational program of \TheoryName{} is no longer at a stage where the completion-state shell-trace core \(\TraceCore\) should be treated as the default unexplained stopping point. The trace-core collapse theorem had already shown that the full completion-core graph presentation \(\CompletionCore\) carries no extra benchmark-facing freedom beyond the smaller trace core \(\TraceCore\). The later ambient trace-core forcing theorem then spent the preferred next move available there: \(\TraceCore\) itself is now forced by the deeper completion-balance ambient dynamics package \(\AmbientPackage\).

The present note records the routing consequence of that theorem. It does not reopen the ambient mathematics, and it does not claim a completed first-principles substrate derivation of gravity. It records only which kind of next manuscript would now count as an honest benchmark-facing gain below the current ambient-dynamics frontier.

\paragraph{Framework and claim boundary.}
\TheoryProgram{} denotes the broader \Aether{} / \Aflow{} framework built on the claims that the \Aether{} is the underlying four-dimensional substrate of reality, the \Aflow{} is its intrinsic ordered motion, observed three-dimensional space is the local experiential slice of that deeper substrate, S-time is the experienced order of change arising from matter, light, and the \Aflow{}, and observed expansion is the three-dimensional appearance of deeper four-dimensional ordered motion. Within that framework, \TheoryName{} denotes the exact-closure theory in which the effective gravitational dynamics are adopted to be exactly Einsteinian with universal matter coupling. In the active sequence, \emph{adoption} means use of that established relativistic dynamics without claiming substrate derivation, while \emph{derivation} is reserved for a first-principles recovery from explicit substrate variables.

The present note stays strictly inside that boundary. It records only which kind of next manuscript would now count as an honest benchmark-facing gain below the ambient-dynamics frontier.

\section{Fixed Inputs}

\begin{definition}[Ambient-frontier fixed inputs]
The criterion recorded here uses the following fixed inputs.
\begin{enumerate}
    \item The benchmark exact-closure package, which fixes the public one-metric GR target of the repository.
    \item The benchmark gatekeeping layer, including the post-bridge status correction and the upstream-compression safeguard that blocks same-output deeper-origin relays from counting as new front-line progress once the benchmark-relevant object is unchanged.
    \item The direct exact-shell-affine-nucleus forcing theorem, which already proves that the exact-shell affine nucleus \(\AffineNucleus\) is forced by the minimal zero-response affine completion core \(\CompletionCore\).
    \item The trace-core collapse theorem, which proves that the full completion-core graph presentation carries no extra benchmark-facing freedom beyond the smaller completion-state shell-trace core \(\TraceCore\).
    \item The ambient trace-core forcing theorem, which proves that the trace core \(\TraceCore\) is itself a canonical and necessary reduction of the completion-balance ambient dynamics package \(\AmbientPackage\).
\end{enumerate}
\end{definition}

\begin{remark}
The present note does not replace those manuscripts. It records the routing consequence of reading them together under the active safeguard against recursive depth drift.
\end{remark}

\section{Why the Ambient Trace-Core Theorem Is the Frontier}

\begin{definition}[Neutral ambient relay]
A \emph{benchmark-neutral ambient relay} is a manuscript that preserves the same benchmark one-metric output, the same ordinary matter route, the same exact-shell affine nucleus consequence, the same trace-core consequence, and the same claim boundary while merely restating the trace core \(\TraceCore\), restating the full completion-core presentation \(\CompletionCore\), restating the nucleus \(\AffineNucleus\), replaying a preserved older ambient-style relay, or relocating the unexplained layer deeper without forcing or shrinking the ambient-balance package itself.
\end{definition}

\begin{proposition}[The ambient trace-core forcing theorem is the live frontier]
At the current stage of the primitive-reservoir observer-reduction line, the theorem deriving the minimal completion-state shell-trace core \(\TraceCore\) from completion-balance ambient dynamics \(\AmbientPackage\) is the live benchmark-facing frontier.
\end{proposition}

\begin{proof}
That theorem changes benchmark-facing status at current scope because it removes the trace core \(\TraceCore\) itself as the live internal stopping point. The current bridge datum is no longer merely a smaller trace-level package; it is now a canonical output of the deeper ambient package \(\AmbientPackage\). By contrast, a manuscript that preserves the same benchmark package, the same trace-core consequence, and the same ambient stopping point while merely rephrasing the trace core, restating the larger completion-core presentation, or replaying an older ambient relay does not alter benchmark-facing status unless it adds a comparably sharp gain. Therefore the live frontier is the ambient trace-core forcing theorem rather than the earlier trace-core collapse theorem or a later same-output relay.
\end{proof}

\begin{corollary}[Status of the trace core as a next-step target]
Under the current routing safeguard, the completion-state shell-trace core \(\TraceCore\) is no longer itself the default next benchmark-facing target.
\end{corollary}

\begin{proof}
The ambient trace-core forcing theorem already proves that \(\TraceCore\) is forced by \(\AmbientPackage\). Once that is on record, another manuscript that spends the move on the trace core again does not remove any further benchmark-relevant freedom. The open burden has therefore moved below the trace core itself.
\end{proof}

\section{The Next Benchmark Criterion}

\begin{definition}[Admissible next benchmark-facing ambient manuscript]
At the current ambient-dynamics frontier, a manuscript counts as an admissible next benchmark-facing move only if it does at least one of the following:
\begin{enumerate}
    \item derives or uniquely forces the completion-balance ambient dynamics package \(\AmbientPackage\) from still more primitive explicit substrate structure in a way that does more than merely relocate the unexplained layer deeper;
    \item proves a genuinely smaller theorem-level collapse, sharper necessity result, or sharper uniqueness result strictly inside \(\AmbientPackage\);
    \item does either of the previous two while preserving the same benchmark one-metric package, the same ordinary matter route, and the same claim boundary.
\end{enumerate}
\end{definition}

\begin{theorem}[Next honest benchmark-facing task at ambient scope]
The next honest benchmark-facing manuscript at the current ambient-dynamics frontier must either derive or uniquely force the completion-balance ambient dynamics package \(\AmbientPackage\) from still more primitive explicit substrate structure in a benchmark-relevant way, or else prove a genuinely smaller collapse, sharper necessity result, or sharper uniqueness result inside that ambient package. A manuscript that only restates the trace core \(\TraceCore\), only restates the full completion-core presentation \(\CompletionCore\), only restates the exact-shell affine nucleus \(\AffineNucleus\), or only repeats the same ambient trace-core consequence without such a new gain does not satisfy the criterion.
\end{theorem}

\begin{proof}
By Proposition 1, the live frontier already lies at the theorem that forces the current trace core from the ambient package. Therefore a second manuscript below it must improve on that status rather than merely accompany it. A deeper-origin manuscript can do so only if it changes benchmark-facing status by more than depth alone. If it simply reproduces the same ambient trace-core consequence through another relay without forcing or shrinking the ambient package itself, the routing rule classifies it as benchmark neutral. The remaining honest alternatives are therefore exactly those stated in the definition: either a more primitive derivation or uniqueness forcing of \(\AmbientPackage\), or a sharper internal collapse, necessity result, or uniqueness result inside \(\AmbientPackage\).
\end{proof}

\section{What the Next Move Should Not Be}

\begin{proposition}[Screened next moves]
At the current ambient-dynamics frontier, none of the following is the default next benchmark-facing move:
\begin{enumerate}
    \item another restatement of the completion-state shell-trace core \(\TraceCore\);
    \item another restatement of the full completion-core graph presentation \(\CompletionCore\);
    \item another restatement of the exact-shell affine nucleus \(\AffineNucleus\);
    \item another replay of preserved ambient-style relays as if they again defined the live burden without forcing or shrinking the ambient package itself;
    \item a manuscript that introduces a second operative metric, enlarges the ordinary matter law, or uses stronger promotion language.
\end{enumerate}
\end{proposition}

\begin{proof}
Items (1)--(4) fail the preceding criterion because they do not directly address the current ambient burden. They revisit either an already-forced smaller bridge object, an already-collapsed larger internal presentation, or the same ambient stopping point rather than removing the remaining ambient-balance freedom or proving a smaller internal datum inside \(\AmbientPackage\). Item (5) violates the fixed claim boundary of the benchmark package and therefore cannot count as a disciplined next step on the active line. The benchmark package remains the exact one-metric GR package already fixed by adoption, so any admissible next manuscript must preserve that boundary.
\end{proof}

\begin{remark}
This screening statement is not a denial that the trace-core theorem, the completion-core theorem, or the older ambient-style theorems matter historically. It is a routing statement: they are not the honest way to spend the next benchmark-facing move from the current ambient-dynamics frontier.
\end{remark}

\section{Conclusion}

The positive result of this note is routing discipline at the current ambient-dynamics frontier. The ambient trace-core forcing theorem remains the live benchmark-facing gain because it already removes the trace core as the internal stopping point on the active line. The trace core therefore no longer sets the honest next burden, and neither do the larger completion-core presentation or the already-forced exact-shell affine nucleus.

The scope remains narrow and honest. The benchmark package is unchanged: the one operative metric is still exactly \(\CompletedMetric\), the ordinary matter route is unchanged, there is no second operative metric, and the low-energy matter law is not enlarged. Nothing here upgrades adoption to derivation or licenses stronger promotion language.

The next open burden is therefore clear. The next benchmark-facing manuscript should derive or uniquely force the completion-balance ambient dynamics package \(\AmbientPackage\) from still more primitive explicit substrate structure in a way that changes benchmark-facing status, unless the sharper honest gain available instead is a genuinely smaller collapse, sharper necessity result, or sharper uniqueness result strictly inside that ambient package. That is the clean next manuscript target at current scope.

\end{document}
